\begin{references}
\item The Women’s Health Initiative Study Group. Design of the women’s health initiative clinical trial and observational study. Control Clin Trials. 1998;19:61-109.
\sourcepage{367}
\item Califf R. Biomarker definitions and their applications. Exp Biol Med. 2018;243:213-221.
\item Black DM, Bauer DC, Vittinghoff, E, et al. Treatment-related changes in bone mineral density as a surrogate biomarker for fracture risk reduction: meta-regression analyses of individual patient data from multiple randomised controlled trials. Lancet Diabetes Endocrinol. 2020;8(8):672-682.
\item Nissen SE, Wolski K. Rosiglitazone revisited: an updated meta-analysis of risk for myocardial infarction and cardiovascular mortality. Arch Intern Med. 2010;170(14):1191-1201.
\item The Action to Control Cardiovascular Risk in Diabetes Study Group. Effects of intensive glucose lowering in type 2 diabetes. N Engl J Med. 2008;358:2545-2559.
\item Barter PJ, Caulfield M, Eriksson M et al. Effects of torcetrapib in patients at high risk for coronary events. N Engl J Med. 2007;357:2109-2122.
\item Self WH, Semler MW, Leither LM, et al. Effect of hydroxychloroquine on clinical status at 14 days in hospitalized patients with COVID-19: a randomized clinical trial. JAMA. 2020;324:2165-2176.
\item Ridker PM, Everett BM, Pradhan A, et al. Low-dose methotrexate for the prevention of atherosclerotic events. N Engl J Med. 2019;380(8):752-762.
\item Friedman LM, Furberg C, DeMets DL. Fundamentals of Clinical Trials. 4th ed. Springer; 2010.
\item Avins AL. Can unequal be more fair? Ethics, subject allocation, and randomised clinical trials. J Med Ethics. 1998;24:401-408.
\item Subak LL, Wing R, West DS, et al. Weight loss to treat urinary incontinence in overweight and obese women. N Engl J Med. 2009;360(5):481-490.
\item Prentice RL, Caan B, Chlebowski RT, et al. Low-fat dietary pattern and risk of invasive breast cancer: the women’s health initiative randomized controlled dietary modification trial. JAMA. 2006;295:629-642.
\item CDP Research Group. The coronary drug project. Initial findings leading to modifications of its research protocol. JAMA. 1970;214:1303-1313.
\item Yu WY, Lu B, Tan D, et al. Effect of topical brimonidine on alcohol-induced flushing in Asian individuals: A randomized clinical trial. JAMA Dermatol. 2020;156(2):182-185.
\item Noseworthy JH, Ebers GC, Vandervoort MK, et al. The impact of blinding on the results of a randomized, placebo‐controlled multiple sclerosis clinical trial. Neurology. 1994;44(1):16.
\item Al-Lamee R, Thompson D, Dehbi HM, et al. Percutaneous coronary intervention in stable angina (ORBITA): a double-blind, randomised controlled trial. Lancet. 2018;391(10115):31-40.
\item Cummings SR. Clinical trials without clinical sites. JAMA Intern Med. 2021;181:680-684.
\item Chesnut CH 3rd, Silverman S, Andriano K, et al. A randomized trial of nasal spray salmon calcitonin in postmenopausal women with established osteoporosis: the prevent recurrence of osteoporotic fractures study. PROOF Study Group. Am J Med. 2000;109(4):267-276.
\item Cummings SR, Chapurlat RD. What PROOF proves about calcitonin and clinical trials. Am J Med. 2000;109(4):330-331.
\item Huang AJ, Cummings SR, Schembri M, et al. Continuous transdermal nitroglycerin therapy for menopausal hot flashes: a single-arm, dose-escalation trial. Menopause. 2016;23(3):330-334.
\item Pfeffer M, Stevenson L. Beta-adrenergic blockers and survival in heart failure. N Engl J Med. 1996;334:1396-1397.
\item Armitage P, McPherson C, Rowe B. Repeated significance tests on accumulating data. J R Stat Soc. 1969;132A:235-244.
\item Parker MJ, Pryor G, Gurusamy K. Hemiarthroplasty versus internal fixation for displaced intracapsular hip fractures: a long-term follow-up of a randomised trial. Injury. 2010;41(4):370-373.
\item McClung MR, Boonen S, Torring O, et al. Effect of denosumab treatment on the risk of fractures in subgroup of women with postmenopausal osteoporosis. J Bone Mineral Res. 2012;27:211-218.
\item Wang R, Lagakos SW, Ware JH, et al. Statistics in medicine—reporting of subgroup analyses in clinical trials. NEJM. 2007;357:2189-2194.
\end{references}
